Uma das principais dificuldades encontradas durante o desenvolvimento do trabalho foi a qualidade das imagens de ultrassom. Para que uma imagem seja considerada de boa qualidade, é necessário seguir um protocolo de aquisição e análise rigoroso, como explicado por \cite{Mercadante2010}. Cada imagem deve conter em posições específicas pontos anatômicos bem definidos, o que não foi possível devido à quantidade limitada de imagens disponíveis, por conta disso, os modelos treinados podem propagar erros que seriam facilmente evitáveis com imagens de melhor qualidade. A presença de artefatos, ruídos e variações na qualidade das imagens pode comprometer a precisão dos modelos de segmentação e, consequentemente, a acurácia das medidas numéricas derivadas dessas segmentações. Outro desafio foi o desenvolvimento da U-Net Fuzzy, com poucos artigos sendo encontrados na literatura sobre o tema, o que dificultou a implementação e a obtenção de resultados satisfatórios. A falta de uma base teórica sólida e de exemplos práticos na literatura exigiu um esforço adicional para adaptar os conceitos de sistemas neuro-fuzzy à arquitetura da U-Net, o que resultou em um modelo com desempenho inferior ao esperado. 
Vale comentar que os modelos estão intimamente ligados ao aparelho de ultrassom utilizado, já que o formato da imagem pode variar significativamente entre diferentes equipamentos, o que pode limitar a generalização dos resultados para outras situações. Uma possível solução para esse problema seria a utilização de técnicas que retirem informações específicas do aparelho, deixando apenas as informações relevantes para a segmentação.

Por conta de como as regras no módulo fuzzy são implementadas, a quantidade de parâmetros aumenta  de forma significativa com o número de regras, o que limita a quantidade de regras que podem ser utilizadas. No caso da U-Net Fuzzy, 9 regras foram utilizadas, o que resultou em um modelo com 8.060.478 parâmetros, enquanto a U-Net multitarefa possui 2.729.249 parâmetros. Ainda sim, as redes desenvolvidas são relativamente pequenas, essa característica representa uma vantagem em termos de eficiência computacional, pois favorece a implementação em ambientes com recursos limitados e abre a possibilidade de adaptação do aplicativo desenvolvido para dispositivos móveis.

Como trabalhos futuros, podem ser estudados e implementados modelos que utilizem o algoritmio de atenção, desenvolvido pela Google \cite{Vaswani2017}, em conjunto com a U-Net, como também a utilização de técnicas de transferência de aprendizado (\textit{fine tuning}), que podem ajudar a melhorar o desempenho dos modelos, especialmente quando a quantidade de dados disponíveis é limitada.  

